% ==============================================================================
% SIAM journal submission template (skeleton)
% Official SIAM classes:
%   * siamltex.cls      (legacy, 1995)  —  historically on CTAN as the `siam`
%                                         package, but old and brittle
%   * siamart220829.cls (current, 2022) —  distributed by SIAM, NOT in TeX Live
% To submit to a SIAM journal, download siamart220829.cls (and the .bst files)
% from https://www.siam.org/publications/journals/about-siam-journals/information-for-authors
% and replace the preamble below with:  \documentclass{siamart220829}
%
% Editorial style (SIAM "Author Submission Guidelines"):
%   Layout     : single column, 11pt, double-spaced for review
%   Body font  : Times Roman
%   Math font  : Times (mathptmx)  /  Computer Modern acceptable
%   Captions   : 10pt, sentence case
%   Citation   : numbered [1]  (natbib numeric)
%   Figures    : single column only (~3.5in wide in the printed article)
%   Sections   : 1. Introduction, 2. Background, ..., N. Conclusion
% References  : siamplain.bst  (numeric, alphabetical by author)
%               siamunbib.bst  (unsorted, citation order)
% ==============================================================================
\documentclass[11pt]{article}
% For the real SIAM class (download siamart220829.cls from siam.org), use:
% \documentclass[review]{siamart220829}
% Or the legacy class (in TeX Live as the `siam` package):
% \documentclass{siamltex}

\usepackage[T1]{fontenc}
\usepackage[utf8]{inputenc}
\usepackage{mathptmx}            % Times text + math (SIAM preference)
\usepackage{helvet}
\usepackage{courier}
\usepackage{amsmath,amssymb,amsfonts}
\usepackage{amsthm}              % SIAM theorem environments
\usepackage{graphicx}
\usepackage{booktabs}
\usepackage{array}
\usepackage{xcolor}
\usepackage{lineno}

% ---- Citations: numbered, natbib ----
\usepackage[numbers,sort&compress]{natbib}
\bibliographystyle{plain}        % close to siamplain (alphabetical, numeric)
% With the real SIAM class, use:
%   \bibliographystyle{siamplain}     % alphabetical, numbered
%   \bibliographystyle{siamunbib}     % unsorted, citation order

\linenumbers                     % review mode: line numbers

% ---- SIAM-style theorem environments ----
\newtheorem{theorem}{Theorem}
\newtheorem{lemma}[theorem]{Lemma}
\newtheorem{proposition}[theorem]{Proposition}
\newtheorem{corollary}[theorem]{Corollary}
\theoremstyle{definition}
\newtheorem{definition}[theorem]{Definition}
\newtheorem{example}[theorem]{Example}
\theoremstyle{remark}
\newtheorem{remark}[theorem]{Remark}

% ---- SIAM-style section numbering ----
\renewcommand{\thesection}{\arabic{section}}
\renewcommand{\thesubsection}{\thesection.\arabic{subsection}}

% \usepackage{hyperref}           % load last if used

\title{A Concise and Informative Title for the SIAM Submission}
\author{First Author\thanks{Department of X, University of Y, City, Country
  (\texttt{first.author@univ.edu}).} \and Second Author\thanks{Department of Z,
  Institute of W, City, Country.}}
\date{}

\begin{document}
\maketitle

% ==============================================================================
% Abstract  (SIAM: unstructured, ~150-250 words, no cites, math allowed)
% ==============================================================================
\begin{abstract}
% TODO: 150-250 words. State the problem, the approach, the main result (with
% a key formula or theorem statement if applicable), and the significance.
% No citations. Inline math ($x$) is fine.
\end{abstract}

% ==============================================================================
% 1. Introduction
% ==============================================================================
\section{Introduction}
\label{sec:intro}
% TODO: motivate the problem; summarise prior limitations; state contributions.
% SIAM style: open with the broad problem, then the specific question, then
% the approach, then a bulleted or enumerated contribution list, then a
% paper roadmap. Cite prior work with \cite{ref_example}.

% ==============================================================================
% 2. Background / Related Work
% ==============================================================================
\section{Background}
\label{sec:background}
% TODO: required preliminaries, notation, and prior results that the rest of
% the paper builds on. State lemmas with \begin{lemma}...\end{lemma}.

\begin{lemma}
\label{lem:example}
% TODO: state a prior result that the analysis will rely on.
\end{lemma}

% ==============================================================================
% 3. Problem Formulation
% ==============================================================================
\section{Problem Formulation}
\label{sec:problem}
% TODO: formal problem statement, assumptions, and the quantity of interest.

\begin{definition}
\label{def:example}
% TODO: define a key term.
\end{definition}

% ==============================================================================
% 4. Main Results
% ==============================================================================
\section{Main Results}
\label{sec:results}
% TODO: state and prove the main theorems; discuss tightness; give corollaries.

\begin{theorem}
\label{thm:main}
% TODO: state the main result precisely.
\end{theorem}

\begin{proof}
% TODO: proof of Theorem~\ref{thm:main}.
\end{proof}

% ==============================================================================
% 5. Numerical Experiments
% ==============================================================================
\section{Numerical Experiments}
\label{sec:experiments}
% TODO: setup (datasets / parameters / baselines), main convergence / error
% tables, runtime comparison, and a discussion of how the numerics corroborate
% the theory.

\begin{figure}[t]
\centering
% \includegraphics[width=0.8\textwidth]{figures/result.pdf}
\fbox{\rule{0pt}{5cm}\rule{0.7\textwidth}{0pt}}
\caption{Figure caption goes BELOW the figure (SIAM convention).
Captions are sentence case, 10pt.}
\label{fig:result}
\end{figure}

\begin{table}[t]
\caption{Table caption goes ABOVE the table (SIAM convention).}
\label{tab:results}
\centering
\begin{tabular}{lcc}
\toprule
Method   & Error          & Time (s)       \\
\midrule
Baseline & $1.2 \times 10^{-2}$          & 12.4          \\
Ours     & $\mathbf{3.4 \times 10^{-4}}$ & $\mathbf{8.1}$ \\
\bottomrule
\end{tabular}
\end{table}

% ==============================================================================
% 6. Conclusion
% ==============================================================================
\section{Conclusion}
\label{sec:conclusion}
% TODO: recap the main theorem, its significance, and open questions.

% ==============================================================================
% References
% Compile:  pdflatex main  ->  bibtex main  ->  pdflatex main  ->  pdflatex main
% ==============================================================================
\bibliography{references}
% Inline fallback (uncomment if not using BibTeX):
% \begin{thebibliography}{99}
% \bibitem{ref_example} F.~Author and S.~Author, \emph{Title of the paper},
%   SIAM J.~X, YY (20YY), pp.~1--25.
% \end{thebibliography}

\end{document}
